\section{Introdução}
\label{sec:intro}

Trinta e cinco milhões de brasileiros --- 35\% da força de trabalho ocupada --- atuam em ocupações com exposição quase nula à inteligência artificial ($\theta \le 2$ em uma escala de 0 a 10). Não é a sofisticação das tarefas que as isola, mas o caráter predominantemente manual e informal de suas atividades. O trabalho doméstico (6{,}7 milhões de postos) e a construção civil estrutural (3{,}2 milhões) reúnem mais empregos do que todas as ocupações de tecnologia, finanças e administração somadas; sua exposição ao índice de \citet{eloundou2024gpts} é de 0{,}1 e 0{,}9, respectivamente, contra 7{,}5 dos desenvolvedores de software e 8{,}6 dos operadores de máquinas de escritório.

Esse quadro qualifica a narrativa dominante sobre IA e mercado de trabalho. Nos Estados Unidos, \citet{eloundou2024gpts} estimam que cerca de 80\% da força de trabalho tem ao menos 10\% de suas tarefas afetadas por grandes modelos de linguagem. No Brasil, a mesma métrica revela um padrão assimétrico: a exposição é \emph{positivamente} associada à escolaridade e à renda, concentrando-se no emprego formal. Em economias em desenvolvimento, o mercado de trabalho conta com um amortecedor estrutural ausente nos países de alta renda --- a informalidade \citep{gmyrek2023generative} ---, e esse amortecedor situa-se exatamente na base do espectro de exposição.

Este artigo investiga quem, no mercado de trabalho brasileiro, está exposto à IA e por meio de quais mecanismos econômicos. Desenvolvo um modelo de designação de tarefas com setor informal em que a IA individual está universalmente disponível, mas a IA corporativa exige capital organizacional restrito à formalidade. Testo as previsões do modelo com os microdados trimestrais da Pesquisa Nacional por Amostra de Domicílios Contínua (PNAD Contínua/IBGE, referente ao primeiro trimestre de 2026, com mais de 227 mil trabalhadores analisados individualmente), combinados ao mapeamento determinístico entre a COD brasileira e a classificação O*NET-SOC.

A análise empírica revela que a exposição à IA ordena-se fortemente com a escolaridade: com controles mincerianos completos (idade, idade ao quadrado, gênero e raça) e erros-padrão agrupados por ocupação (122 grupos), cada ano adicional de estudo associa-se a 0{,}23 ponto a mais de exposição tecnológica (erro-padrão 0{,}03; $p < 0{,}001$; $N = 227.629$). Em paralelo, a exposição exibe relação inversa e substantiva com a informalidade: na relação incondicional, cada ponto de exposição associa-se a $-6{,}23$ p.p. de informalidade ($p < 0{,}001$). Uma vez condicionada à escolaridade individual, o coeficiente da exposição atenua-se em $37{,}2\%$ para $-3{,}91$ p.p. (erro-padrão 0{,}96), enquanto cada ano de estudo associa-se a $-2{,}61$ p.p. de informalidade ($p < 0{,}001$), evidenciando a mediação parcial pelo capital humano aliada à retenção de um canal tecnológico direto. Por fim, o gradiente educacional acentua-se nas Unidades da Federação com maior maturidade do setor formal: a interação entre escolaridade e formalidade da UF (\emph{leave-one-out}) alcança 0{,}28 (erro-padrão 0{,}06; $p < 0{,}001$).

A interpretação desses padrões é de \emph{sorting} ocupacional de equilíbrio, não de impacto causal ou deslocamento imediato de trabalhadores. O modelo formaliza essa distinção: trabalhadores mais escolarizados possuem vantagem comparativa em tarefas cognitivas e abstratas complementares à IA, autosselecionando-se nas ocupações mais expostas do setor formal, onde o retorno ao capital organizacional corporativo compensa os custos de conformidade tributária e regulatória.

Este estudo articula três linhagens de pesquisa. Dialoga, inicialmente, com a literatura que mensura a exposição das ocupações a novas tecnologias \citep{frey2017future,acemoglu2018race,acemoglu2022tasks,eloundou2024gpts}, na qual a automação desloca tarefas existentes mas também cria novas atividades \citep{acemoglu2019automation}, de modo que a tecnologia frequentemente complementa o trabalho humano \citep{autor2015why}. Integra, em seguida, os modelos de escolha contratual entre formalidade e informalidade em economias em desenvolvimento \citep{ulyssea2018firms}. Por fim, apoia-se na teoria de alocação de trabalhadores a tarefas e pareamento no mercado de trabalho \citep{acemoglu2011skills,autordorn2013growth}.

O trabalho avança em três dimensões. Em primeiro lugar, mensura a exposição ocupacional à IA no nível de microdados para o mercado de trabalho brasileiro, incorporando controles demográficos mincerianos completos e inferência agrupada por ocupação. Em segundo lugar, formaliza o prêmio de capital organizacional corporativo como mediador da viabilidade econômica da adoção tecnológica sob informalidade contratual. Por último, orienta o desenho de políticas públicas ao evidenciar que os programas de requalificação e transição tecnológica devem priorizar o estrato médio-formalizado, segmento diretamente exposto à automação informacional.

O restante do artigo organiza-se da seguinte forma. A Seção~\ref{sec:modelo} formaliza o modelo teórico de designação; a Seção~\ref{sec:dados} descreve os microdados e o mapeamento ocupacional; a Seção~\ref{sec:estrategia} detalha a identificação empírica; a Seção~\ref{sec:resultados} discute as estimativas e a bateria de robustez; e a Seção~\ref{sec:conclusao} conclui.
